\section*{Overview}

Prime factorization is the step that exposes arithmetic structure. Once a number is written as a product of prime powers,
questions about divisors, gcd/lcm, multiplicative functions, and modular conditions usually become much clearer.

\section*{When to Use It}

Use factorization when:

\begin{itemize}
  \item divisor counts or divisor sums matter,
  \item the statement talks about prime powers or coprimality,
  \item a multiplicative function needs to be evaluated from the prime exponents,
  \item the input size suggests the factors themselves are more useful than the raw number.
\end{itemize}

\section*{Core Idea}

Every positive integer has a unique factorization:
\[
n = p_1^{e_1} p_2^{e_2} \cdots p_k^{e_k}.
\]
The implementation question is not whether this decomposition exists, but which method finds it fast enough:

\begin{itemize}
  \item trial division,
  \item sieve-assisted factorization,
  \item Pollard Rho for harder 64-bit cases.
\end{itemize}

\section*{Key Insight}

Factorization is usually the preprocessing, not the end goal. The real value is what it unlocks afterward: divisor
enumeration, phi, Mobius, CRT conditions, and multiplicative formulas.

\section*{Operations / Main Technique}

\begin{itemize}
  \item divide by each prime factor while it still fits,
  \item count the exponent of each factor,
  \item if the remaining value is \(> 1\), it is prime in the trial-division setting.
\end{itemize}

\paragraph{Worked example.}
If \(n = 360\), then
\[
360 = 2^3 \cdot 3^2 \cdot 5.
\]
From that, the divisor count is \((3+1)(2+1)(1+1) = 24\).

\section*{Correctness Intuition}

Whenever a prime \(p\) divides \(n\), dividing out all copies of \(p\) records its exponent exactly. After all primes up
to \(\sqrt{n}\) are tested, any remaining value must itself be prime, because a composite remainder would still have a
factor at most its square root.

\section*{Complexity Analysis}

\begin{itemize}
  \item trial division: \(O(\sqrt{n})\),
  \item repeated queries with smallest-prime-factor sieve: roughly \(O(\log n)\) per number after preprocessing,
  \item large 64-bit inputs may need Miller-Rabin plus Pollard Rho.
\end{itemize}

\section*{Implementation}

The code includes:

\begin{itemize}
  \item basic trial-division factorization,
  \item divisor generation from the prime-exponent list.
\end{itemize}

That is enough for many problems before the harder randomized methods become necessary.

\section*{Common Pitfalls}

\begin{itemize}
  \item Trial-dividing up to the original \(\sqrt{n}\) instead of the shrinking current value.
  \item Forgetting the leftover prime factor when the loop ends.
  \item Building divisor formulas from distinct primes but ignoring their exponents.
\end{itemize}

\section*{Variants / Extensions}

\begin{itemize}
  \item SPF sieve for many small-to-medium queries.
  \item Pollard Rho for difficult 64-bit factorization.
  \item Use factorization as the base for phi, Mobius, divisor sums, and multiplicative functions.
\end{itemize}

\section*{Practice Problems}

\begin{itemize}
  \item Count divisors or sum divisors.
  \item Factor many input values after one sieve.
  \item Construct all divisors from the prime powers.
\end{itemize}

\section*{References}

\begin{itemize}
  \item \href{https://cp-algorithms.com/algebra/factorization.html}{cp-algorithms: Integer Factorization}.
  \item \href{https://usaco.guide/gold/prime-factorization}{USACO Guide: Prime Factorization}.
  \item \href{https://codeforces.com/catalog?locale=en}{Codeforces Catalog}.
\end{itemize}
